%% complexity -- the part of the paper that is not mechanized: complexity and
%% sharper bounds. A framed preface; the remark on the complexity of the
%% procedure; the smaller unit reader and the antecedent rejection clock with
%% the sharper bound for (U) (sec:un-smaller); leaf cones and tree classes, and
%% signed satisfiability with copy literals (sec:cx-facts); witness lengths, the
%% summary search and the 2-EXPSPACE / 3-EXPTIME upper bound, arity one
%% (PSPACE) and arity zero; the exponential lower bound on witnesses, the
%% embedding of binary trees into arity n >= 2, and PSPACE-hardness for every
%% fixed arity n >= 2. Proofs omitted here: appendix-complexity.tex (app:cx).

\providecommand{\parity}{\operatorname{par}}
\providecommand{\CondI}{\textup{(I0)}}
\providecommand{\CondU}{\textup{(U)}}
\providecommand{\CondII}{\textup{(II)}}
\providecommand{\GuardS}{\Guard_{\mathrm{s}}}
\providecommand{\Lnm}{\mathsf L_{\mathrm{nm}}}
\providecommand{\RU}{\mathcal R_{\mathrm U}}
\providecommand{\PosOp}{\operatorname{Pos}}
\providecommand{\cxcode}{\operatorname{code}}   % prefix-free binary code of a word
\providecommand{\cxemb}{E}                      % embedding of binary trees into arity n

\section{Complexity and Sharper Bounds (Not Mechanized)}\label{sec:cx}

\noindent\fbox{\parbox{\dimexpr\linewidth-2\fboxsep-2\fboxrule\relax}{%
\emph{This section is not mechanized.} Every result before it, up to and
including the main theorem (\cref{thm:as-main}) and its applications
(\cref{sec:ap}), is mechanized (\cref{sec:mechanization}), and nothing before
it uses a result of this section. The section collects the complexity of the
problem and bounds sharper than those of the main theorem. Its proofs use the
mechanized results freely; those omitted here are in \cref{app:cx}, which is
not mechanized either.}}

\medskip
Throughout, $s=\len X+\len Z$ is the number of variables after flattening
(\cref{sec:bg-constraints}), $v=\len X$, and $E_1(x)=2^x$,
$E_{k+1}(x)=2^{E_k(x)}$. \Cref{sec:un-smaller} gives a smaller reader for
condition \CondU, with a sharper bound, and \cref{sec:cx-facts} records two
facts on tree classes and signed systems. The upper bounds of
\cref{sec:cx-upper} hold for every arity $n$ and every variance $c$, with
constants that do not depend on $n$; the arity enters only the size of the
alphabet and of the input. The lower bounds of \cref{sec:cx-lower} are proved
for a covariant constructor of arity at least two.

\begin{remark}[Complexity]\label{rem:as-complexity}
The procedure of \cref{sec:as-procedure} is written to make its correctness
evident: \cref{as:step-u,as:step-ii} enumerate all words up to a triply
exponential length. The applications of \cref{sec:ap} need only the statement
of the main theorem. This section gives sharper procedures, the complexity of
the problem, and lower bounds.
\end{remark}

\subsection{A Smaller Reader for Condition (U)}\label{sec:un-smaller}

The bound of \cref{thm:un-decidable} is not the smallest one. Two observations
give a much smaller bound. First, most states of the source automaton
(\cref{def:un-source}) can never occur on a rejected word, so they can be
merged into one sink. Second, the depth bound needs only the first rejected
cut, and on the words that the unit reader rejects this cut is determined by
the antecedent alone, which is read by a much smaller automaton. These results
serve the complexity bounds of \cref{sec:cx-upper}; the decision procedure and
its correctness do not use them.

\begin{lemma}[The reduced unit reader]\label{lem:un-reduced}
Let $p\in X^\pm$, $v=\len X$ and $s=\len X+\len Z$. Let $\widehat{\mathcal R}_p$
arise from the reachable part of $\RU(p)$ by replacing every state
$\zeta=(P,b,P')$ with $b=1$ or with $\mathrm{du}(\zeta)$ false by one new
absorbing state $\bot\notin\Rej$: every transition into a replaced state is
redirected to $\bot$, and no pair that involves $\bot$ is admitted. Then:
\begin{enumerate}[leftmargin=*]
\item a word is rejected by $\widehat{\mathcal R}_p$ iff it is rejected by
  $\RU(p)$, and then both readers have the same state at every cut, so they
  have the same holes;
\item $\widehat{\mathcal R}_p$ has the rejected-path property;
\item $\widehat{\mathcal R}_p$ has at most $1+2^{v+1}(2^s-1)<2^{v+s+1}$ states.
\end{enumerate}
No satisfiability assumption is needed.
\end{lemma}

The replaced states are closed under transitions, because the bit $b$ never
returns to $0$ and $\mathrm{du}$, once false, stays false
(\cref{lem:un-source}(2)); so a rejected word never meets them. The count uses
sign coherence: on a kept state, $P'$ is nonempty and fixes the common sign of
all elements of $P$ and $P'$. The proof is in \cref{app:cx}.

\begin{lemma}[Antecedent rejection clock]\label{lem:un-clock}
Let $p\in X^\pm$ and $v=\len X$, and let $\mathcal A_p$ be the reader whose
automaton is the reachable part of the upper automaton of $p$ in $\varphi$,
whose admission relation is empty, and whose rejected states are the sets $P$
with $P\cap\DeadT^\eps_\varphi=\varnothing$. It has at most $2^{v+1}-1$
states. If $w$ is rejected by $\RU(p)$, then $w$ is rejected by $\mathcal A_p$
with the same first rejected cut, and if moreover $w\ne\eps$, then
$\depth(w)\le2(2^{v+1}-1)(r(w)+1)-1$.
\end{lemma}

On a word rejected by $\RU(p)$, the bit and $\mathrm{du}$ are constant, so only
the first component of the source automaton decides rejection; the depth bound
is then \cref{thm:bg-depth} for $\mathcal A_p$, which uses neither admissions
nor holes. The proof is in \cref{app:cx}.

\begin{proposition}[A sharper bound for (U)]\label{prop:un-sharper}
Let $\varphi$ be satisfiable and $X\ne\varnothing$, let $v=\len X$ and
$s=\len X+\len Z$, and put
\[
 A_{\mathrm U}=2^{v+1},\qquad K_{\mathrm U}=2^{v+s+1},\qquad
 D_{\mathrm U}=2A_{\mathrm U}(4K_{\mathrm U}-1)-1 .
\]
If \CondU\ fails, there are $p\in X^\pm$ and a word $w$ with
$\len w\le\holeBound(K_{\mathrm U},D_{\mathrm U})$ such that $(p,w)$ is live in
$\varphi$ and $\DeadU_\psi(p,w)$. Moreover
$\holeBound(K_{\mathrm U},D_{\mathrm U})\le E_3(5v+3s+16)$.
\end{proposition}

\begin{proof}
If \CondU\ fails, some $\RU(p)$ has a hole (\cref{lem:un-holes}), hence so does
$\widehat{\mathcal R}_p$, which has $m\le K_{\mathrm U}$ states and the
rejected-path property (\cref{lem:un-reduced}). By \cref{thm:bg-brt} it has a
hole $w_0$ with $r(w_0)\le4m-2$; if $w_0=\eps$, we are done. Otherwise $w_0$ is
rejected by $\RU(p)$, and \cref{lem:un-clock} gives
$\depth(w_0)\le2(2^{v+1}-1)(4m-1)-1\le D_{\mathrm U}$. \Cref{thm:bg-search}
for $\widehat{\mathcal R}_p$ with $d=D_{\mathrm U}$ gives a hole of length at
most $\holeBound(m,D_{\mathrm U})\le\holeBound(K_{\mathrm U},D_{\mathrm U})$,
since $\holeBound$ is nondecreasing in each argument
(\cref{app:bg-constants}). This is a hole of $\RU(p)$ by
\cref{lem:un-reduced}(1), hence a word as stated by \cref{lem:un-holes}.

For the estimate, write $A=A_{\mathrm U}$, $K=K_{\mathrm U}$ and
$D=D_{\mathrm U}$. Then $D+1\le8AK$ and $K+D\le9AK$, so the quantity $\Xi$ of
\cref{lem:bg-estimate} satisfies
$\Xi(K,D)\le3K^2+288A^2K^3+3K+24AK+4\le322A^2K^3$, and
$5\Xi(K,D)+3\le1613A^2K^3<2^{11}A^2K^3=2^{5v+3s+16}$. By
\cref{lem:bg-estimate}, $\holeBound(K,D)\le2^{2^{5\Xi(K,D)+3}}\le
E_3(5v+3s+16)$.
\end{proof}

\begin{remark}[Singly exponential cases]\label{rem:un-single}
In two cases the witnesses of \cref{prop:un-sharper} are much shorter. Let
$\varphi$ be satisfiable and $p\in X^\pm$.
\begin{itemize}[leftmargin=*]
\item \emph{No self events.} If $\varphi$ has a parity assignment
  (\cref{rem:ch-parity}), in particular if $f$ is covariant, then no pair of
  reachable states of $\Bsrc_p$ is admitted (\cref{rem:sm-parity}). The holes
  of $\RU(p)$ are then exactly its rejected words, and a shortest one visits no
  state of $\widehat{\mathcal R}_p$ twice, so it is shorter than
  $2^{v+s+1}$.
\item \emph{Arity one.} If $n=1$, every pair of cuts $i<j$ of a word is a
  comparison. If $w$ is a hole of $\RU(p)$, then so is its prefix
  $w_J=\wpre w{J(w)}$, because the comparisons of $w_J$ are comparisons of
  $w$ with the same states. By \cref{lem:un-clock}, $w_J$ is rejected by
  $\mathcal A_p$ with the first rejected cut $J(w)$, and no two of its cuts
  $i<j\le J(w)$ have the same state of $\mathcal A_p$: the unary word would
  then repeat the run of $\mathcal A_p$ with period $j-i$, and the cut
  $J(w)-(j-i)$ would already be rejected. Hence $J(w)<2^{v+1}-1$, and \CondU,
  if it fails, fails at a word of length at most $2^{v+1}-2$.
\end{itemize}
\end{remark}

\subsection{Tree Classes and Signed Systems}\label{sec:cx-facts}

Two facts used below concern the tree classes. Like the rest of this section,
neither is mechanized, and no result before this section uses them. The
first, \cref{lem:lf-classes}, describes finite and regular trees through leaf
cones (\cref{sec:lf}); the lower bounds of \cref{sec:cx-lower} use it, and the
decision procedure and the transfers of \cref{sec:tc} do not. The second,
\cref{lem:bg-signed-sat}, extends the satisfiability criteria of
\citet{Zhao26} (\cref{thm:bg-nsse-sat}(2,3)) to copy literals and to all three
tree classes; the upper bounds of \cref{sec:cx-upper} use it, and the decision
procedure decides satisfiability by \cref{prop:sm-decide,thm:bg-nsse} instead.
Both proofs are in \cref{app:cx}.

\begin{lemma}[Cones and tree classes]\label{lem:lf-classes}
For every tree $t$, the positions of $t$ labeled $f$ are exactly the words
outside $L^0_t\cup L^1_t$. Hence $t$ is finite iff
$\Sigma^*\setminus(L^0_t\cup L^1_t)$ is finite, and $t$ is regular iff $L^0_t$
and $L^1_t$ are regular languages.
\end{lemma}

\begin{lemma}[Signed satisfiability]\label{lem:bg-signed-sat}
For every arity $n\ge0$, every variance $c$, every system $\chi$ over $\BFT$,
copy literals included, and every tree class $\mathcal D$,
\[
 \Sat_{\mathcal D,c}(\chi)\iff\Sat_{\mathcal D}(\chi^\pm),
 \qquad\text{and}\qquad
 \Sat_c(\chi)\iff\Sat_{\mathrm{reg},c}(\chi).
\]
\end{lemma}

\subsection{Upper Bounds}\label{sec:cx-upper}

\paragraph{Witness lengths}
When both systems are satisfiable in the tree class considered and the
entailment fails, a counterexample word (\cref{def:as-certificate}) of length
at most $E_3(10v+12s+15)\le E_3(22s+15)$ exists over $\BFT$, and of length at
most $E_3(20v+24s+15)\le E_3(44s+15)$ over $\BF$ with $f$ covariant
(\cref{thm:as-main}); every counterexample word yields a countermodel
(\cref{lem:as-countermodel}). \Cref{prop:as-bounds} and
\cref{tab:as-bounds} bound the kinds separately. Words for \CondI\ and holes
that are not mortal are singly exponential in $s$; for the latter the pump
keeps the purely periodic case polynomial in the sizes of the automata
(\cref{thm:nm-pump}), for every instance of the two-word problem with self
events that satisfies (FH). The words for \CondU\ and the mortal holes need
the triply exponential estimates of compression. All these bounds are existential: they bound a shortest
counterexample, not every one.

\paragraph{Searching without storing the word}
Enumerating all words up to the bound of \cref{thm:as-main} decides the
problem in triply exponential space. The words need not be stored. The
compression of \cref{thm:mo-anchored} only ever uses a dictionary of boundedly
many blocks per level, and a level is described, for the purpose of the
levels above it, by the derived monitor and channel on these blocks, which
the column summaries determine (\cref{lem:mo-summary}). A nondeterministic
search can therefore build the hierarchy from the bottom, one level at a time,
keeping only the current derived tables.

\begin{proposition}[Summary search]\label{prop:cx-search}
Let $\mathcal M_0$ be a lettered monitor with $N\ge1$ states over a finite
alphabet $\Gamma$, let $\Theta_0$ be a positive channel on it, both given by
their tables, and let $d\ge0$. Whether $(\mathcal M_0,\Theta_0)$ has a positive
hole of depth at most $d$ is decidable in nondeterministic space polynomial in
$N+d+1$, $\sGsharp(N,d)$ and $\len\Gamma$, hence in deterministic space
$2^{O((N+d+1)^3)}\cdot\mathrm{poly}(\len\Gamma)$.
\end{proposition}

\begin{proof}
Put $g=\sGsharp(N,d)$. We use the summaries, column words and agreement of
\cref{app:bg-search} and their extensions by channels in \cref{app:mo}. The
search keeps a level $j\le d$ and a \emph{current pair} $(\mathcal W,\Theta)$:
a lettered monitor with the $N$ states of $\mathcal M_0$ and $j$ symbolic
states, and a positive channel on it, over a current alphabet. Initially
$j=0$ and $(\mathcal W,\Theta)=(\mathcal M_0,\Theta_0)$, over $\Gamma$; for
$j\ge1$ the current alphabet is $\{1,\ldots,h\}$ with $h\le g$. The search
repeats one of two moves.
\begin{itemize}[leftmargin=*]
\item \emph{Accept} if some letter $z$ and some $m\le3(N+j)$ make $z^m$ a
  positive hole of $(\mathcal W,\Theta)$. This is checked by computing the
  labels of $z^m$ and testing every pair of cuts.
\item \emph{Derive}, if $j<d$: choose a marker $\ell$ and a number $1\le h\le g$,
  and guess a valid column word with $h$ tracks over the current alphabet
  without $\ell$, one column at a time, while computing the positive summary
  of the tracks read so far by the automaton of \cref{lem:mo-summary}, starting
  from the summary of $h$ empty blocks; a binary counter stops the guess before
  $\Pcol(N+j,h)$ columns. Require that the final equality fields declare
  the $h$ blocks pairwise distinct. Replace $(\mathcal W,\Theta)$ by the
  restriction of $(\mathcal W_\ell,\Theta_\ell)$ to the $h$ blocks, named
  $1,\ldots,h$, which the final summary determines at every state
  (\cref{app:bg-search}, step~4, and \cref{lem:mo-summary}(2)), and increase
  $j$.
\end{itemize}

\emph{Space.} A current pair over $h\le\max(g,\len\Gamma)$ letters with
$M\le N+d+1$ states is given by its transitions, final targets, the relations
$\Lambda$, $\Lfin$, $R$, $S$, $S'$ on labels, the predicate $F$ and the set
$\Phi$, in $O(M^2h^2+Mh\log M)$ bits. A positive summary has at most
$\Pcol(M,h)$ values, hence $O(M^2h^2+Mh\log M)$ bits, and so does the
counter; a column has $h$ entries. Only the current pair, one summary, one
column, the counter and $j$ are stored. Since $g\le2^{\Xisharp(N,d)}$ by
\cref{thm:mo-anchored} and $\Xisharp(N,d)\le30(N+d+1)^3$, the nondeterministic
space is polynomial in $N+d+1$, $g$ and $\len\Gamma$, and Savitch's theorem
squares it, which keeps the stated form.

\emph{Soundness.} We show by induction on the number of moves that there are
markers $\widetilde\ell_0,\ldots,\widetilde\ell_{j-1}$ and an injective map
$\rho_j$ from the current alphabet into the letters of level $j$, with a valid
image family for these markers (\cref{app:bg-search}), such that
$(\mathcal W,\Theta)\equiv_{E,\rho_j}(\mathcal M^{\widetilde{\boldsymbol\ell}}_j,
\Theta^{\widetilde{\boldsymbol\ell}}_j)$, where $E$ is the current alphabet.
Initially $\rho_0$ is the identity. In a derive move, the guessed column word
spells blocks $Y_1,\ldots,Y_h$ over $E$ that avoid $\ell$; by
\cref{lem:mo-summary}(1) the computed summary is theirs, so the new pair is the
restriction of $(\mathcal W_\ell,\Theta_\ell)$ to $\{Y_1,\ldots,Y_h\}$. By
\cref{lem:mo-agreement},
$(\mathcal W_\ell,\Theta_\ell)\equiv_{\{Y_a\},\rho_j^*}
((\mathcal M^{\widetilde{\boldsymbol\ell}}_j)_{\rho_j(\ell)},
(\Theta^{\widetilde{\boldsymbol\ell}}_j)_{\rho_j(\ell)})$, and the images avoid
$\rho_j(\ell)$. Put $\widetilde\ell_j=\rho_j(\ell)$ and
$\rho_{j+1}(a)=\rho_j^*(Y_a)$, which is injective because the $Y_a$ are
distinct and $\rho_j$ is injective. If the search accepts with $z^m$, then
$\rho_j(z)^m$ is a positive hole of $(\mathcal M^{\widetilde{\boldsymbol\ell}}_j,
\Theta^{\widetilde{\boldsymbol\ell}}_j)$, since agreement preserves positive
holes (\cref{app:mo-agreement}). Expanding it along the markers gives, by
\cref{lem:bg-desub,lem:mo-desub}, a positive hole $w$ of
$(\mathcal M_0,\Theta_0)$. Each expansion ends with its marker and its blocks
avoid it, so the canonical hierarchy of $w$ recovers the expanded levels and
reaches the unary word $\rho_j(z)^m$ at level $j$, and $\depth(w)\le j\le d$.

\emph{Completeness.} Suppose a positive hole of depth at most $d$ exists. By
\cref{lem:mo-dictionary} there is one, $w$, of depth $D\le d$ whose
dictionaries satisfy $g_i\le g$ for $1\le i\le D$; let $(w_j,\ell_j)$ be its
hierarchy, $\mathcal A_j$ its dictionaries and $(\mathcal M_j,\Theta_j)$ its
towers. We follow a run in which, for $1\le j\le D$, the current pair is the
restriction of $(\mathcal M_j,\Theta_j)$ to $\mathcal A_j$, up to the names of
the letters. At level $j<D$, choose $\ell=\ell_j$ and $h=g_{j+1}$. The letters
of $\mathcal A_{j+1}$ are $h$ distinct blocks over $\mathcal A_j$ that avoid
$\ell_j$, and their positive summary depends only on the tables at the letters
of $\mathcal A_j$, so it is also their summary for the current pair. By the
deletion argument in the proof of \cref{lem:mo-columns}, some valid column word
with fewer than $\Pcol(N+j,h)$ columns has the same final summary; the
search guesses it, and the new pair is the restriction of
$(\mathcal M_{j+1},\Theta_{j+1})$ to $\mathcal A_{j+1}$, because the summary
determines it. At level $D$, $w_D=z^m$ is a positive hole of
$(\mathcal M_D,\Theta_D)$ by \cref{cor:mo-tower}, hence of the current pair, and
\cref{lem:mo-unary} shortens it to $m\le3(N+D)$. The search accepts.
\end{proof}

The \emph{trivial channel} $(\mathcal Q\times\Gamma,\varnothing,\varnothing,
\varnothing,\varnothing)$ makes every nonempty word positive, so its positive
holes are the holes of the monitor, and \cref{prop:cx-search} also decides
whether a monitor has a nonempty hole of depth at most $d$.

\begin{theorem}[Complexity of the decision]\label{thm:cx-upper}
Existential entailment over $\BFT$, for every arity and variance, and over
$\BF$ with $f$ covariant, is decidable in each of the three tree classes in
deterministic space $E_2(O(s))\cdot\mathrm{poly}(m)$, where $m$ is the size of
the input, with the arity, the variance and the systems given explicitly. It is
therefore in 2-EXPSPACE and in 3-EXPTIME.
\end{theorem}

\begin{proof}
By \cref{lem:bg-flattening}, flattening is linear, so let the systems be flat.
Over $\BF$, apply the reduction of \cref{thm:tc-gadget}, which is polynomial
and doubles $s$ (\cref{cor:tc-bf}). For $n=0$ enumerate the $3^s$ assignments
(\cref{rem:tc-nullary}), in polynomial space. Let $n\ge1$.

\emph{Satisfiability.} For each class $\mathcal D$, $\Sat_{\mathcal D}$ of a
system $\chi$ is decided in polynomial time: by \cref{lem:bg-signed-sat}, it is
satisfiability of the signed system of the system $\chi^\flat$ of
\cref{lem:bg-sandwich}, and by \cref{thm:bg-nsse-sat}(1) this is the
absence of a label clash, and over finite trees also of a cycle clash, which are
closure and graph tests of polynomial size; regular satisfiability is
satisfiability. If $\varphi$ is not $\mathcal D$-satisfiable, the entailment
holds; if $\varphi$ is and $\psi$ is not, it fails. Otherwise both are
satisfiable, and \cref{thm:tc-regular,thm:tc-finite} reduce the problem to
arbitrary trees.

\emph{The three conditions.} By \cref{cor:ex-entailment}, the entailment holds
iff \CondI\ and \CondU\ hold and no extracted instance has a hole. Condition
\CondI\ is decided in polynomial space: the closure $D_\chi$ of
\cref{def:sm-automata} is computed by saturation in polynomial time, so is
every transition of the bound automata, whose states are sets of signed
variables, and the tests in the proof of \cref{prop:sm-decide}, together with
the computation of $\Fup\chi$ (\cref{lem:sm-mu}), are reachability questions
in these automata or in products of two of them, each with one bit, and are
decided in nondeterministic, hence deterministic, polynomial space. For \CondU,
compute the tables of the reduced unit readers $\widehat{\mathcal R}_p$ of
\cref{lem:un-reduced}, with fewer than $K_{\mathrm U}=2^{v+s+1}$ states each,
in space $2^{O(s)}\mathrm{poly}(m)$ (\cref{lem:un-source}). By
\cref{lem:un-holes,lem:un-reduced}, \CondU\ fails iff some
$\widehat{\mathcal R}_p$ has a hole. The empty word is tested directly. As in
the proof of \cref{prop:un-sharper}, by \cref{thm:bg-brt,lem:un-clock}, a
reduced unit reader with a hole has the empty hole or a nonempty hole of depth
at most $D_{\mathrm U}=2^{v+2}(4K_{\mathrm U}-1)-1$, and its nonempty holes are
those of its lettered monitor (\cref{def:bg-monitor}). So
\cref{prop:cx-search} with the trivial channel decides \CondU\ in space
$2^{O((K_{\mathrm U}+D_{\mathrm U}+1)^3)}\mathrm{poly}(n)=E_2(O(s))\,\mathrm{poly}(n)$.

For \CondII, fix $p,q\in X^\pm$. The instance $I_{pq}$ satisfies (FH), so by
\cref{cor:nm-short-holes} it has a hole iff it has a hole with
$\len u+\len t<\Lnm(N_A,N_B,N_{A'},N_{B'})$, a bound singly exponential in $s$ (\cref{prop:as-bounds}), or a mortal hole. The first
alternative is decided by enumerating such pairs and checking the conditions of
\cref{def:ex-ss2self} cut by cut, in space $2^{O(s)}\mathrm{poly}(m)$. For the
second, by \cref{lem:mo-pi-reader,lem:mo-depth}, $I_{pq}$ has a mortal hole iff
$\RPi$ has a $\Pi$-hole of depth at most $d_I$; this reader has
$\Nbar\le2^{2v+4s+2}$ states (proof of \cref{cor:mo-extracted}), so
$d_I<2^{4v+4s+3}$, and its tables, including the mortality of the states of $A$,
are computed by reachability in space $2^{O(s)}\mathrm{poly}(m)$. Its
$\Pi$-holes are the positive holes of its lettered monitor with the base channel
(\cref{lem:mo-base}), and \cref{prop:cx-search} decides their existence in space
$E_2(O(s))\,\mathrm{poly}(n)$.

\emph{Total.} All tests together use deterministic space
$E_2(O(s))\,\mathrm{poly}(m)$, and $s\le m$. A deterministic machine in this space
has at most $2^{E_2(O(s))\mathrm{poly}(m)}$ configurations, which gives the time
bound.
\end{proof}

\begin{remark}[Plain entailment]\label{rem:cx-plain}
Plain entailment $\varphi\models_{\mathcal D}x\le y$ is the case
$\psi=\{x\le y\}$, $Z=\varnothing$. \Cref{thm:cx-upper} therefore also places
plain entailment, for every arity and variance and in each tree class, in
2-EXPSPACE. The earlier procedure for plain entailment places non-entailment in
N3EXPTIME and
entailment in coN3EXPTIME, hence in 3EXPSPACE
\nssecite{Corollary ``Complexity upper bounds''}; the summary search improves
the deterministic space upper bound to 2-EXPSPACE, one exponential lower.
\end{remark}

\paragraph{Arity one}
For $n=1$ there is a single letter, words are determined by their lengths, and
a suffix comparison is a comparison of lengths. Here a direct shortening gives
singly exponential witnesses for every instance, mortal or not.

\begin{theorem}[Arity one]\label{thm:cx-unary}
Let $n=1$, with either variance.
\begin{enumerate}[leftmargin=*]
\item Let $I$ be an instance of the two-word problem with self events that
  satisfies (FH), and let $r_I\le N_AN_{A'}$ be the number of states of
  $A\times A'$ reachable from the start. If $I$ has a hole, it has a hole
  $(u,t)$ with $\len t\le N_B-1$ and $\len u+1+\len t\le r_I+2N_B-2$.
\item Every reader with $N$ states over a one-letter alphabet that has a hole
  has a hole of length at most $N-1$.
\item Let $\varphi$ and $\psi$ be flat systems over $\BFT$ with $s$ variables.
  If they are $\mathcal D$-satisfiable and
  $\varphi\not\models_{\mathcal D}\Exists\psi$, a counterexample word of length
  less than $2^{4s+3}$ exists. Over $\BF$ with $f$ covariant, if $\varphi$ and
  $\psi$ are $\mathcal D$-satisfiable over $\BF$ and the entailment fails, the
  pair $(\varphi^\natural,\psi^\natural)$ of \cref{thm:tc-gadget}, which has $2s$
  variables, has a counterexample word of length less than $2^{8s+3}$.
  Existential entailment at arity one is in PSPACE, for both variances and in
  each tree class, over $\BFT$ and, for covariant $f$, over $\BF$.
\end{enumerate}
\end{theorem}

\begin{proof}
Write the letter as $0$. For words $0^a$ and $0^b$ and cuts $x\le a$, $y\le b$,
the pair $(x,y)$ is a cross comparison of $(0^a,0^b)$ iff $b-y\le a-x$, and it is
aligned iff $b-y=a-x$; every pair of cuts $i<j$ of a word is a comparison.

(1) Let $(u,t)$ be a hole, $J=J(t)$, and let $(x,e)$ be a head event; by (FH),
$e\le J$ (\cref{lem:ex-head-before-J}). Since $\len u-x=\len t-e\ge\len t-J$,
we may delete $\len t-J$ letters from the end of both words: let $u'$ be the
prefix of $u$ of length $L=\len u-(\len t-J)$ and $t'$ the prefix of $t$ of
length $J$. The cuts of $u'$ and $t'$ are cuts of $u$ and $t$ with the same
states, and $J(t')=J$. For $x\le L$ and $y\le J$, the pair $(x,y)$ is a cross
comparison of $(u',t')$ iff $J-y\le L-x$, iff $\len t-y\le\len u-x$, iff it is
one of $(u,t)$; so $(u',t')$ has no event, $t'$ has no self event, and the head
event $(x,e)$ is a head event of $(u',t')$, with $x\le L$. The spine $u'$ is
safe. So $(u',t')$ is a hole whose target is first rejected at its end.

The states $B_0(t'),\ldots,B_J(t')$ are pairwise distinct: if $B_i=B_j$ with
$i<j\le J$, the run of the single letter repeats with period $j-i$ from $i$ on,
so $B_{J-(j-i)}(t')=B_J(t')\in\Rej$ with $J-(j-i)<J$, contradicting the choice
of $J$. Hence $\len{t'}=J\le N_B-1$.

If $L<J$, then $\len{u'}+1+\len{t'}\le2J\le2N_B-2$. Otherwise let $m$ be the
least cut of $u'$ at which $A\times A'$ is in the state it has at the cut
$L-J$. The states of $A\times A'$ at the cuts $0,\ldots,m$ are pairwise
distinct, so $m\le r_I-1$: if two of them coincided, the run would be periodic from the first of them, and the state at the cut $L-J$ would occur before $m$, contrary to the choice of $m$. Let $u''=0^{m+J}$. Its states
at the cuts $x\le m$ are those of $u'$, and at the cut $m+y$, $y\le J$, those of
$u'$ at $L-J+y$. The pair $(u'',t')$ is a hole. Its spine is safe, since all its
states occur on $u'$. The head event $(x,e)$ has $x=L-J+e$ and becomes the
aligned pair $(m+e,e)$, with the same states. A cross comparison $(x,y)$ of
$(u'',t')$ with $x\ge m$ corresponds to the cross comparison $(L-J+x-m,y)$ of
$(u',t')$, with the same states, and one with $x<m$ has the same states as the
pair $(x,y)$ of $(u',t')$, which is a cross comparison because
$L-x>L-m\ge J\ge J-y$; neither is an event. The target and its self events are
unchanged. Finally $\len{u''}+1+\len{t'}\le r_I-1+2J+1\le r_I+2N_B-2$.

(2) Let $w$ be a hole with first rejected cut $J$. Its prefix of length $J$ is
rejected, and its comparisons are pairs of cuts $i<j\le J$, which are
comparisons of $w$ with the same states; so it is a hole. Its states at the cuts
$0,\ldots,J$ are pairwise distinct, as in~(1), so $J\le N-1$.

(3) By \cref{prop:as-bounds}(1), a failure of \CondI\ has a word of
length less than $2^{2v+2s+2}$. A failure of \CondU\ has a \CondU-word of
length at most $2^{v+1}-2$ (\cref{rem:un-single}). A hole of an extracted
instance has one with $\len u+1+\len t\le r_I+2N_B-2$ by~(1), where $r_I\le2^{2v+2s}$ and
$N_B=2^{2v+2s+1}$. Each bound is less than $2^{4s+3}$.
Over $\BF$, the pair of \cref{thm:tc-gadget} is $\mathcal D$-satisfiable over
$\BFT$ and its entailment fails (proof of \cref{thm:as-main}), and it has $2s$
variables, which gives $2^{8s+3}$.
For membership in PSPACE, follow the proof of \cref{thm:cx-upper}, which handles
satisfiability, the tree classes, $\BF$ and \CondI\ in polynomial space, and
replace the searches for \CondU\ and \CondII\ by an enumeration of the
candidates up to these bounds. A unary word is stored by its length, in $O(s)$
bits. The state of an automaton at a cut is a set of signed variables,
recomputed from the start in polynomial space, and each condition of
\cref{def:bg-hole,def:ex-ss2self} quantifies over cuts, which binary counters
enumerate; comparisons are tests of lengths. So all tests use polynomial space.
For $\BF$ the gadget doubles $s$, which keeps the bounds singly exponential.
\end{proof}

\begin{remark}[Arity zero]\label{rem:cx-nullary}
For $n=0$ every tree is $\Bot$, $f()$ or $\Top$ (\cref{rem:bg-nullary}), and the
entailment says that every assignment of $X$ that satisfies $\varphi$ has an
assignment of $Z$ that satisfies $\psi$ (\cref{rem:tc-nullary}). Assignments
have polynomial size and are checked in linear time, so the problem is in
$\Pi^p_2$, hence in PSPACE, in each tree class, which coincide.
\end{remark}

\subsection{Lower Bounds}\label{sec:cx-lower}

\paragraph{Long witnesses}
Shortest counterexamples can be exponentially long, already for plain
entailment over $\BF$ with a binary covariant constructor and antecedents
without literals $u=\Bot$. The construction uses a small automaton for the set
of words that are not extensions of one exponentially long word, the record of
a binary counter.

\begin{proposition}[Exponential witnesses]\label{prop:cx-length}
For every $k\ge1$ there are a word $W_k\in\{0,1\}^*$ of length
$L_k=2(k+1)2^k$ and a satisfiable flat system $\varphi_k$ over $\BF$, with $f$
binary and covariant, without hidden variables and without literals $u=\Bot$,
over $9k+30$ variables, such that $\varphi_k\models x\le y$ fails over arbitrary
and over regular trees, and:
\begin{enumerate}[leftmargin=*]
\item every word of $\Kpos(\rho x)\setminus\Kpos(\rho y)$, for every solution
  $\rho$ of $\varphi_k$, has the prefix $W_k$, and some regular solution has
  $W_k\in\Kpos(\rho x)\setminus\Kpos(\rho y)$;
\item for the instance of \cref{thm:tc-gadget} over $\BFT$, conditions \CondI\
  and \CondU\ hold, and every hole $(u,t)$ of an extracted instance has
  $\len t\ge L_k$; so every counterexample word of that instance
  (\cref{def:as-certificate}) is a hole word longer than $L_k$.
\end{enumerate}
For every arity $n\ge3$, statements (1) and (2) hold for a system over $\BF$
with $f$ covariant of arity $n$ and $18k+61$ variables, which may contain
literals $u=\Bot$.
Hence shortest counterexample words can have length $2^{\Omega(s)}$, for every
fixed arity $n\ge2$.
\end{proposition}

\begin{proof}
\emph{A counter.} Over the alphabet $\Gamma=\{0,1,\#,\$\}$ let
$H=C_0\#C_1\#\cdots\#C_{2^k-1}\$$, where $C_j$ is the binary expansion of $j$
with $k$ bits, least significant bit first; $\len H=(k+1)2^k$. Consider the
nondeterministic automaton $\mathcal N$ over $\Gamma$, all of whose states are
final, with the following states and transitions; a transition that is not
listed does not exist.
\begin{itemize}[leftmargin=*]
\item $\mathsf E$ loops on every letter.
\item $\mathsf F_0,\ldots,\mathsf F_k$ check the format: $\mathsf F_i$ goes on a
  bit to $\mathsf F_{i+1}$ if $i<k$ and to $\mathsf E$ if $i=k$; on $\#$ to
  $\mathsf F_0$ if $i=k$ and to $\mathsf E$ otherwise; on $\$$ to $\mathsf E$ if
  $i<k$.
\item $\mathsf I$ checks that $C_0$ is zero: it goes on $0$ to $\mathsf I$ and
  on $1$ to $\mathsf E$.
\item $\mathsf C(c)$, $c\in\{0,1\}$, tracks the carry of an increment and
  spawns checkers: on a bit $b$ it goes to both $\mathsf C(c\land b)$ and
  $\mathsf W(b\oplus c,k)$, and on $\#$ to $\mathsf C(1)$.
\item $\mathsf W(b,j)$, $0\le j\le k$, waits for the corresponding bit of the
  next block: for $j>0$ it goes on every letter other than $\$$ to
  $\mathsf W(b,j-1)$; $\mathsf W(b,0)$ goes on the bit $1-b$ to $\mathsf E$.
\item $\mathsf O(o)$, $o\in\{0,1\}$, records whether the current block is all
  ones so far: on a bit $b$ it goes to $\mathsf O(o\land b)$; on $\#$ to
  $\mathsf E$ if $o=1$ and to $\mathsf O(1)$ if $o=0$; $\mathsf O(0)$ goes on
  $\$$ to $\mathsf E$.
\end{itemize}
The initial states are $\mathsf F_0$, $\mathsf I$, $\mathsf C(1)$ and
$\mathsf O(1)$, and there are $3k+9$ states. We claim that the words with a run
of $\mathcal N$ are exactly the words in $\Gamma^*\setminus H\Gamma^*$.

No branch reaches $\mathsf E$ on a prefix of $H$: the format is right; $C_0$ is
zero; a checker spawned at the bit of offset $p$ in block $C_i$ reads the $k$
letters between it and the bit of offset $p$ in $C_{i+1}$, which is
$b\oplus c$ for the bit $b$ and the incoming carry $c$, because $C_{i+1}=C_i+1$;
a $\#$ follows only blocks that are not all ones; and $\$$ follows the all-ones
block $C_{2^k-1}$. After $H$ no state remains: $\mathsf F_k$, $\mathsf C(\cdot)$,
$\mathsf O(1)$ and the checkers have no transition on $\$$, and $\mathsf I$ has
none on $\#$. So no word in $H\Gamma^*$ has a run. Every proper prefix of $H$
has a run through the format states. Let $w$ first differ from $H$ at
position $i<\len H$. If $H[i]$ is a bit of $C_j$ and $w[i]\in\{\#,\$\}$, the
format branch, at some $\mathsf F_p$ with $p<k$, reaches $\mathsf E$. If $w[i]$
is the other bit, then $\mathsf I$ reaches $\mathsf E$ if $j=0$, and otherwise
the checker spawned at the same offset of $C_{j-1}$ is at $\mathsf W(H[i],0)$
and reaches $\mathsf E$. If $H[i]=\#$, a bit leads $\mathsf F_k$ to $\mathsf E$,
and $\$$ leads $\mathsf O(0)$ to $\mathsf E$, since $C_j$ is not all ones. If
$H[i]=\$$, a bit leads $\mathsf F_k$ to $\mathsf E$ and $\#$ leads $\mathsf O(1)$
to $\mathsf E$. Since $\mathsf E$ loops, $w$ has a run, which proves the claim.

\emph{Binary words.} Code the tokens by $0\mapsto00$, $1\mapsto01$,
$\#\mapsto10$, $\$\mapsto11$, and let $W_k$ be the code of $H$, of length $L_k$.
Let $\mathcal L$ have the states $q$, $(q,0)$ and $(q,1)$ for every state $q$
of $\mathcal N$, all final, the initial states of $\mathcal N$, a transition on
the bit $b$ from $q$ to $(q,b)$, and a transition on the bit $b'$ from $(q,b)$
to every successor of $q$ on the token coded by $bb'$. It has $9k+27$ states.
By induction on the number of complete codes, the base states reachable on the
code of a token word are the states reachable on that word in $\mathcal N$, and
after one further bit the reachable set is nonempty iff that set is. So a
binary word has a run of $\mathcal L$ iff its longest prefix of even length is
the code of a token word with a run of $\mathcal N$, that is, iff it does not
begin with $W_k$. The language of $\mathcal L$ is therefore
$\{0,1\}^*\setminus W_k\{0,1\}^*$.

\emph{The system.} Let $\varphi_k$ have the variables $x$, $y$, $d$ and $Y_q$
for the states $q$ of $\mathcal L$, and the literals
\[
 f(d,d)\le Y_q,\qquad f(Y_r,d)\le Y_q\ \ (q\xrightarrow0r),\qquad
 f(d,Y_r)\le Y_q\ \ (q\xrightarrow1r),\qquad Y_q\le y\ \ (q\text{ initial}).
\]
In every solution, $\Kpos(Y_q)$ contains every word that has a run of
$\mathcal L$ from $q$, by induction on its length, using
\cref{lem:bg-bf-order}; hence $\Kpos(y)\supseteq\{0,1\}^*\setminus W_k\{0,1\}^*$,
and every word of $\Kpos(x)\setminus\Kpos(y)$ begins with $W_k$. Conversely, let
$\Kpos(Y_q)$ be the set of words with a run from $q$, $d=\Bot$, $\Kpos(y)$ the
language of $\mathcal L$, and $\Kpos(x)$ the set of prefixes of $W_k$. These
sets are regular and prefix closed, since all states are final, and they
satisfy every literal, so this is a regular solution with
$W_k\in\Kpos(x)\setminus\Kpos(y)$. This proves~(1); the antecedent is
satisfiable and the entailment fails over arbitrary and regular trees
(\cref{thm:tc-regular}).

(2) The consequent is $\psi=\{x\le y\}$ and $Z=\varnothing$, so $\psi^\natural=
\psi$ has no facts and no forbidden roots: \CondI\ and \CondU\ hold.
Its pure moves are the copies $(x^0,w)\to(y^0,w)$ and $(y^1,w)\to(x^1,w)$. Let
$(u,t)$ be a hole of an extracted instance $I_{pq}$ of $\varphi^\natural$ and
$\psi$. By \cref{thm:ex-reduction}, $(p,t)\to^*(q,u)$ by pure moves of $\psi$,
and the least solution $\rho$ of $\varphi^\natural$ that contains $(p,t)$ omits
$(q,u)$; so $u=t$ and $(p,q)$ is $(x^0,y^0)$ or $(y^1,x^1)$. By
\cref{lem:tc-gadget}(1), $\rho x$ and $\rho y$ are $\Top$-free, so their
positive cones are empty, and $p=y^1$, $q=x^1$. For $\Top$-free covariant trees
the negative cone is the complement of the constructor positions
(\cref{rem:ch-one-sign}), so $t\in\Kpos(\rho x)\setminus\Kpos(\rho y)$. The
restriction of $\rho$ to $X$ solves $\varphi_k$ over $\BF$, and~(1) gives
$\len t\ge L_k$.

For $n\ge3$, apply \cref{prop:cx-embedding} below: the solutions of the
translated antecedent are, on the original variables, the images of the
solutions of $\varphi_k$, with the same sets of constructor positions, and the
translation has $2(9k+30)+1$ variables; it contains the literal $o=\Bot$. Statement~(2) carries over by the same
argument. Finally $s\le18k+61$ and $\log_2L_k\ge k+1$, so
$L_k\ge2^{(s-61)/18}$.
\end{proof}

The antecedents have no finite solutions, since $\Kpos(y)$ is infinite, so the
proposition gives no lower bound over finite trees.

\paragraph{Larger arities}
A lower bound for a binary constructor does not transfer to a constructor of
larger arity by \emph{bare padding}, that is, by filling the new arguments of
every constructor with $\Bot$. The translated antecedent constrains values that
are not translations of binary trees, and such values may refute the
translated entailment.

\begin{example}[Bare padding fails]\label{ex:cx-padding}
Over $\BF$ with $f_2$ binary and covariant,
$\mathrm{true}\models\exists a\,b\ \ x\le f_2(a,b)$ holds in every tree class:
take the children of $x$, or $a=b=\Bot$ if $x=\Bot$. With $f_3$ ternary and the
new argument filled by $\Bot$, the entailment
$\mathrm{true}\models\exists a\,b\ \ x\le f_3(a,b,\Bot)$ fails at the finite tree
$x=f_3(\Bot,\Bot,f_3(\Bot,\Bot,\Bot))$. In the same way, padding the antecedent
of \cref{prop:cx-length} leaves $x$ unconstrained, and a value of $x$ with a
constructor at the one-letter word $2$, which lies outside $\Kpos(y)$, is a
countermodel with a witness of length one.
\end{example}

The repair restricts every variable to the image of the binary trees, by a
variant of the gadget of \cref{sec:tc-gadget}.

\begin{proposition}[Embedding binary trees]\label{prop:cx-embedding}
Let $n\ge2$, let $f$ be covariant of arity $n$, and let $f_2$ be binary and
covariant. For a tree $t$ over $\{\Bot,f_2\}$ let $\cxemb(t)$ be the tree over
$\BF$ with $\Kpos(\cxemb t)=\Kpos(t)$. For a system $\chi$ over
$\{\Bot,f_2\}$, let $\chi^{\cxemb}$ replace every literal $f_2(a_0,a_1)\le u$ by
$f(a_0,a_1,o,\ldots,o)\le u$ and every literal $u\le f_2(a_0,a_1)$ by
$u\le f(a_0,a_1,o,\ldots,o)$, for a fresh variable $o$. For $\varphi$ over $X$
and $\psi$ over $X\cup Z$, put
\[
 \varphi^{(n)}=\varphi^{\cxemb}\cup\{o=\Bot\}\cup\bigcup_{x\in X}B(x,m_x),
 \qquad
 \psi^{(n)}=\psi^{\cxemb}\cup\bigcup_{z\in Z}B(z,m_z),
\]
where $B(v,m)=\{v\le m,\ m\le f(m,m,o,\ldots,o)\}$,
with the visible variables $X\cup\{o\}\cup\{m_x\}_{x\in X}$ and the hidden
variables $Z\cup\{m_z\}_{z\in Z}$. Then:
\begin{enumerate}[leftmargin=*]
\item $\cxemb$ is a bijection from the trees over $\{\Bot,f_2\}$ onto the trees
  over $\BF$ whose constructor positions lie in $\{0,1\}^*$; it preserves and
  reflects the order, maps each tree class onto its part of the image, and
  $\cxemb(f_2(a,b))=f(\cxemb a,\cxemb b,\Bot,\ldots,\Bot)$;
\item for every tree class $\mathcal D$,
  $\varphi\models_{\mathcal D}\Exists\psi$ over $\{\Bot,f_2\}$ iff
  $\varphi^{(n)}\models_{\mathcal D}\exists\bar z\,\bar m\,\psi^{(n)}$ over
  $\BF$, and the restrictions to $X$ of the solutions of $\varphi^{(n)}$ are
  exactly the images under $\cxemb$ of the solutions of $\varphi$.
\end{enumerate}
The translation has $2s+1$ variables and size linear in the input for fixed $n$.
\end{proposition}

\begin{proof}
(1) Prefix-closed subsets of $\{0,1\}^*$ are exactly the sets of constructor
positions of binary trees and of the trees in the image, and the order on both
sides is inclusion of these sets (\cref{lem:bg-bf-order}). A tree over $\BF$ is finite iff this set is finite, and regular iff it is a
regular language (\cref{lem:lf-classes}), on both sides.

(2) Let $m$ be a tree over $\BF$ with $m\le f(m,m,\Bot,\ldots,\Bot)$. Then
$\Kpos(m)\subseteq\{\eps\}\cup0\Kpos(m)\cup1\Kpos(m)$, and by induction on length
every word of $\Kpos(m)$ lies in $\{0,1\}^*$. Conversely, the regular tree $\Theta$
with $\Kpos(\Theta)=\{0,1\}^*$ and the finite trees $\Theta_h$ with
$\Kpos(\Theta_h)=\{0,1\}^{<h}$ satisfy this literal, and every tree in the image
of $\cxemb$ lies below $\Theta$, and below $\Theta_h$ if it is finite with all
positions shorter than $h$. So, with $o=\Bot$, a value of $v$ admits a value of
$m$ with $B(v,m)$ in the class $\mathcal D$ iff it lies in the image of
$\cxemb$. On this image, a literal of $\chi^{\cxemb}$ with $o=\Bot$ holds for
$\cxemb\circ\sigma$ iff the corresponding literal of $\chi$ holds for $\sigma$,
by~(1). Both claims of~(2) now follow as in the proof of \cref{thm:tc-gadget},
with $\cxemb$ in place of the identity on $\Top$-free trees.
\end{proof}

\paragraph{Hardness}
\Citet[Theorems~1 and~2]{HengleinRehof98} proved PSPACE-hardness of plain
entailment over $\BFT$ for a binary covariant constructor. We prove
PSPACE-hardness for every fixed arity $n\ge2$, in each tree class, by explicit
reductions from the universality and the inclusion problems of
nondeterministic automata over $\{0,1\}$, which are PSPACE-complete
\citep{StockmeyerMeyer73}. For a word $w\in\{0,1\}^*$ let
$\cxcode(w)=e(w)1$, where $e$ replaces each letter $0$ by $00$ and each $1$ by
$01$. The code is prefix free: if $\cxcode(w')\pre\cxcode(w)$ and
$\len{w'}<\len w$, then position $2\len{w'}$ carries the final $1$ of
$\cxcode(w')$ but the letter $0$ in $\cxcode(w)$. From an automaton for $L$ one
obtains in linear time an automaton for $\cxcode(L)$: replace each transition
by two, the first labeled $0$, and add a transition labeled $1$ from every final
state to one new final state.

\begin{theorem}[PSPACE-hardness]\label{thm:cx-hard}
Let $n\ge2$ be fixed and $f$ covariant. In each of the three tree classes:
\begin{enumerate}[leftmargin=*]
\item existential entailment over $\BF$ is PSPACE-hard;
\item plain entailment over $\BFT$ is PSPACE-hard;
\item consequently existential entailment over $\BFT$ is PSPACE-hard, and so is
  the problem with the variance as part of the input, already on covariant
  instances.
\end{enumerate}
\end{theorem}

We claim no hardness for a fixed variance with a contravariant argument, and
none for $n=1$, where the problem is in PSPACE by \cref{thm:cx-unary}. The
reductions use the extra arguments of $f$ either for an unconstrained filler,
whose positions every condition ignores, or for $\Bot$, under literals that
already confine the relevant values to the letters $0$ and $1$; neither is
bare padding.

\begin{proof}
Let $\Sigma=\{0,\ldots,n-1\}$. For an automaton $\mathcal N$ over $\{0,1\}$
with states $Q$, initial state $q_0$ and final states $F$, write $L_q$ for the
words accepted from $q$.

(1) Let $\mathcal F(\mathcal N)$ have the variables $x$, $d$ and $X_q$ for
$q\in Q$, and the literals
\[
 x\le X_{q_0},\quad X_q\le f(X_{q'},d,\ldots,d)\ (q\xrightarrow0q'),\quad
 X_q\le f(d,X_{q'},d,\ldots,d)\ (q\xrightarrow1q'),\quad X_q=\Bot\ (q\in F).
\]
We claim: a tree $t$ over $\BF$ in a class $\mathcal D$ is the value of $x$ in
some $\mathcal D$-solution iff no word of $L_{q_0}$ is a constructor position
of $t$. If $\pi\in\Kpos(t)$ and $q_0\xrightarrow\pi q$, then by induction on
$\pi$ the subtree of $t$ at $\pi$ lies below $X_q$: an upper literal at a
constructor node passes to the child in the direction of the transition. For
$\pi\in L_{q_0}$ with $q\in F$ this contradicts $X_q=\Bot$. Conversely, let $h$
bound the lengths of the positions of $t$ if $t$ is finite, and put
$h=\infty$ otherwise, with $\Sigma^{\le\infty}=\Sigma^*$. Let $\Kpos(d)=\Sigma^{\le h}$ and let $\Kpos(X_q)$ be the
words of $\Sigma^{\le h}$ with no prefix in $L_q$. These sets are prefix
closed, finite if $h<\infty$ and regular otherwise. If $0w\in\Kpos(X_q)$ and
$q\xrightarrow0q'$, then no prefix of $w$ lies in $L_{q'}$, so
$w\in\Kpos(X_{q'})$, and every $iw\in\Kpos(X_q)$ with $i\ne0$ has
$w\in\Kpos(d)$; so the upper literals hold by \cref{lem:bg-bf-order}, and
symmetrically for the letter $1$. For $q\in F$, $\eps\in L_q$, so $X_q=\Bot$.
Finally $\Kpos(t)\subseteq\Kpos(X_{q_0})$, since $\Kpos(t)$ is prefix closed and
avoids $L_{q_0}$.

Given automata $\mathcal N_1,\mathcal N_2$, let $\mathcal N'_i$ accept
$\cxcode(L(\mathcal N_i))$. Let $\varphi=\mathcal F(\mathcal N'_2)$, with all
its variables visible, and let $\psi=\mathcal F(\mathcal N'_1)$ with fresh
variables except $x$, all of them hidden. Both systems are satisfiable, by the
assignment $\Bot$. By the claim, $\varphi\models_{\mathcal D}\Exists\psi$ iff
every $\Top$-free $t\in\mathcal D$ whose constructor positions avoid
$\cxcode(L(\mathcal N_2))$ also avoids $\cxcode(L(\mathcal N_1))$. If
$L(\mathcal N_1)\subseteq L(\mathcal N_2)$, this holds. If $w\in L(\mathcal
N_1)\setminus L(\mathcal N_2)$, let $t$ be the finite tree whose constructor
positions are the prefixes of $\cxcode(w)$; a prefix of $\cxcode(w)$ in
$\cxcode(L(\mathcal N_2))$ would equal $\cxcode(w)$ by prefix freeness, so $t$
is a countermodel. Inclusion of automata is PSPACE-hard, and the reduction is
polynomial.

(2) Let $\mathcal N$ accept $\cxcode(L(\mathcal N_0))$ for a given automaton
$\mathcal N_0$, and let $C=(00\mid01)^*1$ be the set of all codes and $P$ the set
of proper prefixes of codes, the words $e(w)$ and $e(w)0$. Let $\chi$ have the
variables $x$, $d$, $b$, $r$, $r'$ and $X_q$ for $q\in Q$, and the literals
\[
\begin{gathered}
 x\le X_{q_0},\quad X_q\le f(X_{q'},d,b,\ldots,b)\ (q\xrightarrow0q'),\quad
 X_q\le f(d,X_{q'},b,\ldots,b)\ (q\xrightarrow1q'),\quad X_q=\Bot\ (q\in F),\\
 b=\Bot,\qquad f(r',b,\ldots,b)\le r,\qquad f(r,r,b,\ldots,b)\le r' .
\end{gathered}
\]
We show that $\chi\models_{\mathcal D}x\le r$ iff $L(\mathcal N_0)=\{0,1\}^*$.

Suppose $\mathcal N_0$ is universal, so that $\mathcal N$ accepts every code,
and let $\rho$ be a solution. \emph{The values of $r$.} Call $r$ the state of
the words $e(w)$ and $r'$ that of the words $e(w)0$. By induction along $P$, if
$\pi\in P$ is a position of $\rho r$, then its subtree lies above the value of
its state: from $r$, the literal $f(r',b,\ldots,b)\le r$ gives a node $f$ or
$\Top$ whose child $0$, if it exists, lies above $r'$, and from $r'$ the other
literal gives children $0$ and $1$ above $r$. In particular $\rho r$ is not
$\Bot$ at any position in $P$. \emph{The values of $x$.} Every $\pi\in P$ is a
proper prefix of an accepted code, so an accepting run reaches, on $\pi$, a
state $q$ with an outgoing transition; as in~(1), if $\pi$ is a position of
$\rho x$, its subtree lies below $X_q$, which lies below a constructor tree, so
$\rho x$ is not $\Top$ at $\pi$, and the children of that subtree in the
directions $i\ge2$ are $\Bot$. Every code $c$ is accepted, so if $c$ is a
position of $\rho x$, its subtree lies below $X_q=\Bot$ for a final $q$. The
children of a word of $P$ are words of $P$, codes, and words $\pi i$ with
$\pi\in P$ and $i\ge2$, and $\rho x$ is $\Bot$ at the last two kinds; so every
position of $\rho x$ is of one of these three kinds. Hence at every common
position either $\rho x$ is $\Bot$, or the position lies in $P$, where $\rho x$
has a label in $\{\Bot,f\}$ and $\rho r$ one in $\{f,\Top\}$. So
$\rho x\le\rho r$.

Suppose $w\notin L(\mathcal N_0)$, and let $c=\cxcode(w)$, so that no prefix of
$c$ is accepted by $\mathcal N$, by prefix freeness. Let $x$ be the finite tree
whose constructor positions are the prefixes of $c$, $c$ included. For $q\in Q$
let $X_q$ be the join, the union of constructor positions, of the subtrees of
$x$ at the prefixes $\pi$ of $c$ with $q_0\xrightarrow\pi q$; each is the path
tree along the rest of $c$. Then $x\le X_{q_0}$, $X_q=\Bot$ for $q\in F$, and
the upper literals hold with $b=\Bot$ and $d$ the finite tree with
$\Kpos(d)=\{0,1\}^{\le\len c}$, as in~(1). Let $r_0=r'_0=\Top$,
$r_{h+1}=f(r'_h,\Bot,\ldots,\Bot)$ and $r'_{h+1}=f(r_h,r_h,\Bot,\ldots,\Bot)$. By
induction $r_{h+1}\le r_h$ and $r'_{h+1}\le r'_h$, so $(r_h,r'_h)$ satisfies the
two lower literals, and for $h>\len c$ the tree $r_h$ has the label $\Bot$ at
$c$, where $x$ has $f$. This finite solution refutes $x\le r$.

Universality is PSPACE-hard and the reduction is polynomial; since PSPACE is
closed under complement, this proves~(2). Item~(3) follows from~(1) by
\cref{thm:tc-gadget}, or from~(2), since plain entailment is the case
$Z=\varnothing$ and the covariant instances have the variance $c=0$.
\end{proof}

The reduction of~(2) reproves, for $n=2$, the theorem of
\citet{HengleinRehof98}, in all three tree classes, and extends it to every
fixed arity at least two. Together with \cref{thm:cx-upper,thm:cx-unary}, the
problem is PSPACE-hard and in 2-EXPSPACE for every fixed $n\ge2$ and covariant
$f$, and in PSPACE for $n\le1$; whether it is PSPACE-hard for $n=1$ is open.
